\documentclass[10pt,a4paper]{amsart}
\usepackage[margin=19mm]{geometry}
\usepackage{amsmath,amssymb,mathtools}
\usepackage[scr=rsfs,cal=euler]{mathalfa}
\usepackage{xfrac}
\usepackage[unicode,hidelinks,bookmarks=false]{hyperref}
\newcommand{\ie}{i.e.\ }
\newcommand{\cf}{cf.\ }
\newcommand{\resp}{respectively\ }
\newcommand{\Subsection}{\subsection}
\providecommand{\nobreakdash}{\nobreak\hbox{-}}
\setlength{\emergencystretch}{2em}
\multlinegap0pt
\usepackage{bm}
\usepackage{amssymb}
\newtheorem{prop}{Proposition}
\def\G{\mathscr{G}}
\def\R{\mathbb{R}}
\def\T{\mathbb{T}}
\title{Martingale Projections and Quantum Decoherence}
\author{Lane P. Hughston, Levent A. Meng\"ut\"urk}
\date{Source: arXiv:2509.19491v3 (CC BY 4.0)}
\begin{document}
\maketitle

\noindent\emph{Source and licence.} Martingale Projections and Quantum Decoherence. Version arXiv:2509.19491v3 (CC BY 4.0), distributed by arXiv under the Creative Commons Attribution 4.0 International licence, http://creativecommons.org/licenses/by/4.0/, as recorded in the arXiv licence metadata for this exact version and shown on the abstract page https://arxiv.org/abs/2509.19491v3. Full licence notice: \texttt{licenses/arxiv-2509.19491-CC-BY-4.0.txt}.\par\smallskip

\noindent\textit{Editorial scope.} Counted unit: the article's prop mainproposition (source0.tex lines 420--431). The article's complete original proof of it is delivered with it (source0.tex lines 433--474, one \texttt{proof} environment of 42 lines). No mathematics is written, repaired or re-derived: the statement and the proof are the article's own lines, in the article's own notation, and no retained line is substituted or re-flowed.  Retained uncounted as the source-local context the proof needs: lines 368--372; lines 379--381; lines 387--389; lines 411--413.  Nothing was omitted from any retained span, so the package carries no omission record.  No retained line contains a citation, so the wrapper carries no bibliography and no bibliography entry.  No retained line contains a figure environment, an image-inclusion command or drawing code, so no image is delivered and the package declares no asset dependency.  Editorial formatting only, in addition: an AMS \texttt{amsart} preamble that restates the article's own declarations and macros used by the retained lines, namely the environments \texttt{prop} on separate counters, together with the standard packages those lines need.  The retained environments print in this document as Proposition 1 in order of appearance, whereas the article prints the same items with the numbering of its own full document.  The complete native source of the article as distributed in the arXiv e-print for this exact version is bundled unchanged as \texttt{sources/185612/source0.tex}.
\begin{align}
\mathbb{E}^{\nu^{[t_{0},t_{k-1}]}}\left[\Psi_{t_{k}}^{(i,j)}\right] &= \mathbb{E}^{\nu}\left[ \,\, \langle E^{(i)} |\psi_{t_{k}}\rangle \langle \psi_{t_{k}}| E^{(j)} \rangle \,\, \left| \,\, \pi_{t_{0}}^{(q)}, \ldots, \pi_{t_{k-1}}^{(q)} \,: \, \forall q\in\mathcal{Q} \, \right.\right] \label{conditionalquantumexpressionfirst} \\
&= \mathbb{E}^{\nu}\left[ \,\,  \sqrt{\pi_{t_{k}}^{(i)}\pi_{t_{k}}^{(j)}}\exp\left( \text{i}\left( \theta^{(i)} - \theta^{(j)} \right) \right) \,\, \left| \pi_{t_{0}}^{(q)}, \ldots, \pi_{t_{k-1}}^{(q)} \,: \, \forall q\in\mathcal{Q} \, \right.\right] \notag \\
&= \exp\left( \text{i}\left( \theta^{(i)} - \theta^{(j)} \right) \right)\mathbb{E}^{\nu}\left[ \,\,  \sqrt{\pi_{t_{k}}^{(i)}\pi_{t_{k}}^{(j)}} \,\, \left| \,\, \left|\psi_{t_{0}}\right\rangle, \ldots, \left|\psi_{t_{k-1}}\right\rangle \right.\right] \label{conditionalquantumexpression}
\end{align}

\begin{align}
\left|\mathbb{E}^{\nu^{[t_{0},t_{k}]}}\left[\Psi_{t_{k}}^{(i,j)}\right] \right| &= \sqrt{\pi_{t_{k}}^{(i)}\pi_{t_{k}}^{(j)}} \label{expectationremoved} 
\end{align}

\begin{align}
\left|\psi_{t_{k}}\right\rangle = \boldsymbol{\Gamma}^{(t_0,t_{k-1})\mapsto(t_k)}\left( \left|\boldsymbol{\psi}_{\T_{[t_0,t_{k-1}]}}\right\rangle \right)  = \sum_{q\in\mathcal{Q}} \sqrt{\pi_{t_{k}}^{(q)}} \exp\left( \text{i} \theta^{(q)} \right) \left|E^{(q)}\right\rangle, \hspace{0.1in} \text{$\forall t_{k}\in\T\setminus\{t_0\}$}, \label{rewritetransformation}
\end{align}

\begin{align}
\G_{\T^*} \, \triangleq \, \bigvee_{q\in\mathcal{Q}}\G^{\boldsymbol{\pi}^{(q)}}_{\T^*} \label{wholesigmaalgebra}, 
\end{align}

\begin{prop}
\label{mainproposition}
Let $\left|\psi_{t_{k}}\right\rangle$ admit a path transformation of the form in {\rm (\ref{rewritetransformation})} for a fixed $t_k\in\T\setminus\{t_0\}$. If there exists a projection $\boldsymbol{\Gamma}_{\pi^{(q)}}^{(t_0,t_{k-1})\mapsto(t_k)}: \boldsymbol{\Lambda}(\R_+ \,, \T) \rightarrow \boldsymbol{\Lambda}(\R_+ \,, \T)$ where 
\begin{align}
\pi^{(q)}_{t_{k}} = \boldsymbol{\Gamma}_{\pi^{(q)}}^{(t_0,t_{k-1})\mapsto(t_k)}\left(  \boldsymbol{\pi}_{\T_{[t_0,t_{k-1}]}}^{(q)} \right)  \label{projectionforpi}
\end{align}
and $\boldsymbol{\Gamma}_{\pi^{(q)}}^{(t_0,t_{k-1})\mapsto(t_k)} \in \mathbb{H}^{(-)}_{\nu,\G}(\boldsymbol{\Lambda}(\R_+ \,, \T))$ for every $q\in\mathcal{Q}$, then
\begin{align}
\left| \mathbb{E}^{\nu^{[t_{0},t_{k-1}]}}\left[\Psi_{t_{k}}^{(i,j)}\right] \right| < \left| \Psi_{t_{k-1}}^{(i,j)} \right| \label{decoheremceoffdiagnal}
\end{align}
for every $i,j\in\mathcal{Q}$ where $i\neq j$ at that $t_k\in\T\setminus\{t_0\}$.
\end{prop}

\begin{proof}
Since (\ref{projectionforpi}) holds and 
\begin{align}
\boldsymbol{\Gamma}_{\pi^{(q)}}^{(t_0,t_{k-1})\mapsto(t_k)} \in \mathbb{H}^{(-)}_{\nu,\G}(\boldsymbol{\Lambda}(\R_+ \,, \T))  
\end{align}
for every $q\in\mathcal{Q}$, we have
\begin{align}
\mathbb{E}^{\nu}\left[\boldsymbol{\Gamma}_{\pi^{(q)}}^{(t_0,t_{k-1})\mapsto(t_k)}\left(  \boldsymbol{\pi}_{\T_{[t_0,t_{k-1}]}}^{(q)} \right)  \,\, \left| \,\, \G_{\T_{[t_0,t_{k-1}]}} \right.\right] \leq \pi^{(q)}_{\sup(\T_{[t_0,t_{k-1}]})},\end{align}
which implies that
\begin{align}
\mathbb{E}^{\nu}\left[\pi^{(q)}_{t_{k}}  \,\, \left| \,\, \G_{\T_{[t_0,t_{k-1}]}} \right.\right] \leq \pi^{(q)}_{t_{k-1}}  \label{supermartingalepercforquantum}
\end{align}
for every $q\in\mathcal{Q}$ at that $t_k\in\T\setminus\{t_0\}$. Now denote
\begin{align}
\left|\left| \sqrt{\pi^{(q)}_{t_{k}} }\right|\right|^{\nu^{[t_{0},t_{k-1}]}}_p = \left( \mathbb{E}^{\nu}\left[\left(\sqrt{\pi^{(q)}_{t_{k}}}\right)^p \,\, \left| \,\, \G_{\T_{[t_0,t_{k-1}]}} \right.\right] \right)^{\frac{1}{p}} \hspace{0.1in} \text{for $p\geq 1$}. 
\end{align}
By use of (\ref{supermartingalepercforquantum}) and the monotonicity of $ x \mapsto \sqrt{x}$ we deduce that
\begin{align}
\left|\left| \sqrt{\pi^{(q)}_{t_{k}} }\right|\right|^{\nu^{[t_{0},t_{k-1}]}}_2 \leq \sqrt{\pi^{(q)}_{t_{k-1}}}  \label{pnormexpression}
\end{align}
for every $q\in\mathcal{Q}$ at that $t_k\in\T\setminus\{t_0\}$. Then, using the fact that $\pi_{t_{k}}^{(q)}\geq 0$ for all $q\in\mathcal{Q}$, employing (\ref{conditionalquantumexpressionfirst})-(\ref{conditionalquantumexpression}) and using (\ref{wholesigmaalgebra}), we can write
\begin{align}
\left|\mathbb{E}^{\nu^{[t_{0},t_{k-1}]}}\left[\Psi_{t_{k}}^{(i,j)}\right] \right| &= \mathbb{E}^{\nu}\left[ \,\,  \sqrt{\pi_{t_{k}}^{(i)}\pi_{t_{k}}^{(j)}} \,\, \left| \G_{\T_{[t_0,t_{k-1}]}} \right.\right]. \label{absolutedensitymatrixconditional}
\end{align}
Since $\pi_{t_{k}}^{(q)}$ are linearly independent for $q\in\mathcal{Q}$, by use of conditional H\"older's inequality, we have
\begin{align}
\mathbb{E}^{\nu}\left[ \,  \sqrt{\pi_{t_{k}}^{(i)}\pi_{t_{k}}^{(j)}} \,\, \left| \G_{\T_{[t_0,t_{k-1}]}} \right.\right] < \left|\left| \sqrt{\pi^{(i)}_{t_{k}} }\right|\right|^{\nu^{[t_{0},t_{k-1}]}}_2 \left|\left| \sqrt{\pi^{(j)}_{t_{k}} }\right|\right|^{\nu^{[t_{0},t_{k-1}]}}_2 
\end{align}
for $i,j\in\mathcal{Q}$, $i\neq j$, which, by use of (\ref{pnormexpression}), further implies
 \begin{align}
\mathbb{E}^{\nu}\left[ \,  \sqrt{\pi_{t_{k}}^{(i)}\pi_{t_{k}}^{(j)}} \,\, \left| \G_{\T_{[t_0,t_{k-1}]}} \right.\right] < \sqrt{\pi^{(i)}_{t_{k-1}} \pi^{(j)}_{t_{k-1}}} 
= \left|\mathbb{E}^{\nu^{[t_{0},t_{k-1}]}}\left[\Psi_{t_{k-1}}^{(i,j)}\right] \right| \label{finalinequality}
\end{align}
%
for $i,j\in\mathcal{Q}$, $i\neq j$, where the last equality follows from (\ref{expectationremoved}). Therefore, by use of (\ref{absolutedensitymatrixconditional}) and (\ref{finalinequality}), we see that
%
\begin{align}
\left|\mathbb{E}^{\nu^{[t_{0},t_{k-1}]}}\left[\Psi_{t_{k}}^{(i,j)}\right] \right| < \left|\mathbb{E}^{\nu^{[t_{0},t_{k-1}]}}\left[\Psi_{t_{k-1}}^{(i,j)}\right] \right| = \left| \Psi_{t_{k-1}}^{(i,j)} \right| 
\end{align}
%
holds for every $i,j\in\mathcal{Q}$ where $i\neq j$ at that $t_k\in\T\setminus\{t_0\}$.
\end{proof}

\end{document}
